\section{Gold, Jupiter, \& Wine}

Gornahoor quotes Evola here\footnote{\url{http://www.gornahoor.net/?p=2377}}, regarding the possibility of theurgy and true astrology:

\begin{quotex}
In his \emph{The Hermetic Tradition}, Julius Evola suggests that what he did for alchemy, needs to be done also for astrology and theurgy. This post is a possible beginning for whoever would like to take on that task.
\end{quotex}


One possibility is the use of Marcilio Ficino's writings, including his invaluable letters. Ficino was an ordained Catholic priest, living in Florence during the times of Savaranola, before the Reformation. He could not possibly be expected to be plain \& forthright in the same manner in which Evola could afford to be, or Steiner. As Gornahoor has emphasized, \emph{practice} is paramount, and so most esoteric teachers (including St. Paul) have not bothered to be completely perspicacious; add Ficino's context, \& many of his comments and thoughts become powerfully tendentious.

Ficino believed himself to be born under the sign of Saturn, yet what he meant by this (and he wrote against the literal astrologers) was not that he was doomed to one Fate or the other, but that Saturn would resist him should he attempt to live a complacent life, and would aid him should he choose to follow contemplation and scholarship.

Anyone who has listened to Gustav Holst's \emph{The Planets} (particularly Jupiter)\footnote{\url{http://www.youtube.com/watch?v=Nz0b4STz1lo}} should appreciate the Pythagorean connections Ficino makes between the seven important planets, the seven notes of the chromatic scale, and the seven colors of the rainbow. Ficino believed that Creation had implanted correspondences or seminal influences in the universe which were not only real, not only discernible, not only operative on a rational scale, but in final goal (and perhaps also the art of working) responded and resonated with the human spirit. Thus, there is a connection between the element gold, the planet Jupiter, and the liquid wine. In the animal kingdom, the lion would correspond, or manifest the same spiritual influence, which was revealing itself at various levels and in various kingdoms. Among trees, one could choose the oak. The proper arrangement of such could generate a resonance of the ``cosmic Lyre'', of which man was the crux:

\begin{quotex}
Now fiery things aid the attractive power, earthy things the retentive power, airy the digestive, and watery the expulsive... imitate so far as is possible the action of the heavens. But if you can pass through larger spaces in your motion, you will thereby imitate the heavens all the more, \& will get into contact with more the strength of the celestials, which are diffused everywhere.
\end{quotex}


Ficino's influence here was enormous, because flowing out from the Renaissance, the science of good magic (as opposed to profane, or black, magic) bloomed all over Europe, though separating itself from ``divine matters''. With Ficino, it was all one stream of knowledge. Ficino distinguished within ``idolatry'' between false worship of (say) the literal sun, and the worship of the super-celestial Sun through the regard paid to the influence of the literal one. This distinction could be made because

\begin{quotex}
The parts and members of this huge creature the World, I mean all the bodies that are in it, do in good neighborhood as it were, lend and borrow each others Nature; for by reason that they are linked in one common bond, therefore they have love in common; and by force of this common love, there is among them a common attraction, or tilling of one of them to another. And this indeed is Magick\footnote{\url{http://homepages.tscnet.com/omard1/jportat5.html}}.
\end{quotex}


This doctrine is the logical development of Neo-Platonism as read by Ficino; it is implied once Augustine places the Ideas in the mind of God.

Ficino invoked \emph{Cosmos} in the old Orphic hymn:

\begin{quotex}
O heavens, creator of all things, unconquerable part of the Cosmos, of ancient and venerable origin, beginning and end of all things. O Cosmos, father who illuminates the earth's sphere with its spherical movement, the home of the blessed, turning in a revolving motion. Cosmos, heavenly and earthly, guardian and custodian of all things, encircling the whole, holding in your breast the invincible necessity of Nature, red-eyed, untamed, many-formed, diversifying the universe, all-powerful father of Time, most excellent, blessed daemon. Cosmos, hear our prayers, and grant a quiet life to a pious youth.
\end{quotex}


This invocation at the noon hour, with the music of the lyre, was followed by the making of his fortune, as he received Medici's letter guaranteeing his pension. It enabled him to translate (among other things) Plotinus, Plato, and the \emph{Corpus Hermeticum} (which was brought from Macedonia by an unknown monk to the court of the Medici).

By arranging certain things together, including music, work, \& the elements, one could induce spiritual states which were favorable to Fortune's favorable winds. This was not ``magic'' as commonly understood, but rather, ``understanding the times'' in a very personal sense. Rather than abrogating free will in any way, it would indicate the manner in which possibilities might be encouraged to unfold.

Or again, as Dante had it, ``here, Fate and Love are one'' -- the force that rules the Stars moves the heart.

\begin{quotex}
``Far greater than the heavens is he who made you, Lorenzo, \& you yourself will be greater than the heavens as soon as you resolve upon the task. For those celestial bodies are not to be sought by us outside in some other place; for the heavens in their entirety are within us, in whom the light of life and the origin of heaven dwell.''
\end{quotex}


Ficino believed that the work of Plotinus, Iamblichus, Plato, and others were given by God to help elucidate, secretly, the seeds of the Gospel mystery contained in Scripture. Because God had made the world one way, and not another, and because it had to be so by reason and His nature, man's reason for existence was to re-unite Godhead and that over which dominion was given.

How did this transpire?

\begin{quotex}
For the spirit is an intermediary between the gross body of the world and its soul. For whether the world's body or mundane things have their being directly from the world's soul (as Plotinus \& Porphyry think) or whether the world's body just like its soul has its being directly from God, as is the opinion of our theologians and perhaps Timaeus the Pythagorean, the world does wholly live and breathe, and we are permitted to absorb its spirit.
\end{quotex}


The gross is ruled by the subtle. Or, if you wish,

\begin{quotex}
Plotinus \& Avicenna then show that whatever terrestial things come about on earth arise from the general concurrence of all the higher and lower causes at the same time.
\end{quotex}


Thus, true Science is possible, and idolatry is impossible, only when super-celestial things are primarily known.

Several Academy members, including Pico, died in suspicious circumstances during Ficino's lifetime. Ficino walked a very tight rope, but though others had denied the possibility of even symbolically understanding ``the stars'', Ficino's last work \emph{The Book of the Sun} re-affirmed such. He was, in effect, another Plato without a suspicion towards art, and learned in Christian theology.

\begin{quotex}
Celestial figures by their own motions dispose themselves for acting; for by their harmonious rays and motion penetrating everything, they daily influence our spirit secretly just as overpowering music does openly...
\end{quotex}


Nevertheless, the free will of man in the last instance presided over even these; small wonder, then, that it was no idolatry to ``join'' with them in the act of \emph{adaequatio} or knowing, in which super-celestial good was perceived only through the literal figures, but without ignorant idolatry.

\begin{quotex}
It is increased through the same power through which it was begun: it makes progress through the same power through which it was increased; and finally, it is perfected through the same power through which it made progress. hence our soul will sometime be able to become in a sense all things, and even to become God.
\end{quotex}



\flrightit{Posted on 2011-06-11 by Logres}

